%% notes/theory_generalized_root_datum.tex
%% Self-contained section content: Generalized Root Data of CY3 Threefolds
%% Intended for eventual inclusion in Part III (chapters/theory/)
%% Uses theorem environments and macros from main.tex preamble

\chapter{Generalized Root Data of Calabi--Yau Threefolds}
\label{ch:generalized-root-datum}

% AP43: G(X) is not yet constructed for general CY3. The root datum
% R(X) is well-defined (it is extracted from the geometry of X), but the
% quantum vertex chiral group G(X) with root datum R(X) is the TARGET
% of the programme, not a defined object. The root datum axiomatization
% below is independent of the existence of G(X).
The combinatorics of a Calabi--Yau threefold $X$ (its lattice, intersection form, BPS spectrum) constitute a \emph{generalized root datum} $\mathcal{R}(X)$ in the sense of Borcherds.  The programme of this work (Conjecture~CY-A$_3$) posits that $\mathcal{R}(X)$ is the root datum of a quantum vertex chiral group $G(X)$; the root datum itself is a purely geometric/combinatorial object that exists unconditionally.  This chapter axiomatizes the notion of a generalized CY root datum, identifies the axioms distinguishing CY root data from general Borcherds--Kac--Moody (BKM) root data, defines morphisms and operations, and establishes the fibration construction by which a CY$_2$ root datum and an elliptic curve produce a CY$_3$ root datum via the Borcherds multiplicative lift.

\section{Abstract root data of indefinite type}
\label{sec:abstract-root-data}

We follow Borcherds \cite{Borcherds1988, Borcherds1998} and Gritsenko--Nikulin \cite{GN1996, GN2002}, specializing to the CY setting after the general framework.

\begin{definition}[Generalized root datum]
\label{def:generalized-root-datum}
A \emph{generalized root datum} is a tuple
\[
  \mathcal{R} = (\Lambda, (\cdot, \cdot), \Delta, \mathrm{mult}, W, \rho)
\]
consisting of:
\begin{enumerate}[label=(\roman*)]
  \item A free $\mathbb{Z}$-module $\Lambda$ of finite rank $r$, equipped with a symmetric bilinear form $(\cdot, \cdot) \colon \Lambda \times \Lambda \to \mathbb{Z}$ of indefinite signature $(p, q)$ with $p, q \geq 1$.

  \item A \emph{root system} $\Delta \subset \Lambda \setminus \{0\}$ with a decomposition into three disjoint subsets:
  \[
    \Delta = \Delta^{\mathrm{re}} \sqcup \Delta^{\mathrm{im}}_0 \sqcup \Delta^{\mathrm{im}}_1,
  \]
  the \emph{real roots}, \emph{even imaginary roots}, and \emph{odd imaginary roots}, respectively.

  \item A \emph{multiplicity function} $\mathrm{mult} \colon \Delta \to \mathbb{Z}_{\neq 0}$ satisfying:
  \begin{itemize}
    \item $\mathrm{mult}(\alpha) = 1$ for all $\alpha \in \Delta^{\mathrm{re}}$;
    \item $\mathrm{mult}(\alpha) > 0$ for all $\alpha \in \Delta^{\mathrm{im}}_0$;
    \item $\mathrm{mult}(\alpha) < 0$ for all $\alpha \in \Delta^{\mathrm{im}}_1$.
  \end{itemize}

  \item A \emph{Weyl group} $W \subset \mathrm{O}(\Lambda)$ generated by reflections
  \[
    s_\alpha(v) = v - \frac{2(\alpha, v)}{(\alpha, \alpha)} \alpha, \qquad \alpha \in \Delta^{\mathrm{re}},
  \]
  such that $W$ acts on $\Delta$ preserving the decomposition and the multiplicity function.

  \item A \emph{Weyl vector} $\rho \in \Lambda \otimes \mathbb{Q}$ satisfying the normalization
  \[
    (\rho, \alpha_i) = -\tfrac{1}{2}(\alpha_i, \alpha_i)
  \]
  for every simple real root $\alpha_i \in \Pi^{\mathrm{re}} \subset \Delta^{\mathrm{re}}$.
\end{enumerate}
\end{definition}

\begin{remark}[Relation to Borcherds algebras]
\label{rem:borcherds-algebras}
A generalized root datum $\mathcal{R}$ determines a generalized BKM Lie superalgebra $\mathfrak{g}_\mathcal{R}$ with Cartan subalgebra $\mathfrak{h} = \Lambda \otimes_\mathbb{Z} \mathbb{R}$, even generators $\{e_\alpha, f_\alpha\}_{\alpha \in \Delta^{\mathrm{re}} \sqcup \Delta^{\mathrm{im}}_0}$, and odd generators $\{e_\alpha, f_\alpha\}_{\alpha \in \Delta^{\mathrm{im}}_1}$.  The imaginary roots with $\mathrm{mult}(\alpha) > 0$ contribute $|\mathrm{mult}(\alpha)|$ copies of bosonic generators; those with $\mathrm{mult}(\alpha) < 0$ contribute $|\mathrm{mult}(\alpha)|$ copies of fermionic generators.  The ``super'' refers to this $\mathbb{Z}/2$-grading.
\end{remark}

\subsection{Simple roots and the Gram matrix}
\label{ssec:simple-roots-gram}

\begin{definition}[Simple root system]
\label{def:simple-roots}
A \emph{simple root system} for $\mathcal{R}$ is a subset $\Pi = \Pi^{\mathrm{re}} \sqcup \Pi^{\mathrm{im}} \subset \Delta$ such that:
\begin{enumerate}[label=(\alph*)]
  \item $\Pi^{\mathrm{re}} \subset \Delta^{\mathrm{re}}$ is the set of \emph{real simple roots}, forming a basis for the real root subsystem in the sense that every $\alpha \in \Delta^{\mathrm{re}}$ is $W$-conjugate to some $\alpha_i \in \Pi^{\mathrm{re}}$.
  \item $\Pi^{\mathrm{im}} \subset \Delta^{\mathrm{im}}_0 \sqcup \Delta^{\mathrm{im}}_1$ is the set of \emph{imaginary simple roots}: those imaginary roots not expressible as a sum $\beta + \gamma$ with $\beta, \gamma \in \Delta_+$.
  \item The simple roots span a strictly convex cone: $\sum_i n_i \alpha_i = 0$ with $n_i \geq 0$ implies $n_i = 0$ for all $i$.
\end{enumerate}
\end{definition}

\begin{definition}[Generalized Cartan matrix / Gram matrix]
\label{def:gram-matrix}
The \emph{Gram matrix} of a simple root system $\Pi = \{\alpha_1, \ldots, \alpha_n\}$ is the $n \times n$ symmetric matrix
\[
  A_{ij} = (\alpha_i, \alpha_j).
\]
This is a \emph{generalized Cartan matrix of Borcherds type}: $A_{ii} \in 2\mathbb{Z}_{>0}$ for $\alpha_i \in \Pi^{\mathrm{re}}$, $A_{ii} \leq 0$ for $\alpha_i \in \Pi^{\mathrm{im}}$, and $A_{ij} \leq 0$ for $i \neq j$ (with the integrality condition $2A_{ij}/A_{ii} \in \mathbb{Z}_{\leq 0}$ when $A_{ii} > 0$).
\end{definition}

\subsection{Positivity and the fundamental domain}
\label{ssec:positivity}

\begin{definition}[Positive cone and fundamental polyhedron]
\label{def:positive-cone}
The \emph{positive cone} is
\[
  C(\Lambda)_+ = \{v \in \Lambda \otimes \mathbb{R} \mid (v, v) > 0\} / \mathbb{R}_{>0},
\]
a connected component of the space of positive-norm vectors modulo scaling (which is contractible for signature $(p,q)$ with $p \geq 2$, and has two components when $p = 1$; we choose one).  The \emph{fundamental polyhedron} $\mathcal{P}$ of $W$ is a fundamental domain for the $W$-action on the Tits cone
\[
  \mathcal{T} = \{v \in \Lambda \otimes \mathbb{R} \mid (v, \alpha_i) \leq 0 \text{ for finitely many } \alpha_i \in \Pi^{\mathrm{re}}\},
\]
with faces given by the hyperplanes $\alpha_i^\perp$ for $\alpha_i \in \Pi^{\mathrm{re}}$.
\end{definition}

\subsection{The denominator identity}
\label{ssec:denominator-identity}

\begin{proposition}[Weyl--Kac--Borcherds denominator identity]
\label{prop:wkb-denominator}
\ClaimStatusProvedElsewhere
Let $\mathcal{R}$ be a generalized root datum with Weyl vector $\rho$.  The BKM Lie superalgebra $\mathfrak{g}_\mathcal{R}$ satisfies the denominator identity
\[
  e^{-\rho} \prod_{\alpha \in \Delta_+} (1 - e^{-\alpha})^{\mathrm{mult}(\alpha)}
  = \sum_{w \in W} \det(w) \, w \biggl( e^{-\rho} \sum_{\substack{S \subseteq \Pi^{\mathrm{im}} \\ \text{pairwise orthogonal}}} (-1)^{|S|} \, e^{-\sum_{\alpha \in S} \alpha} \biggr)
\]
where the inner sum runs over all subsets $S$ of pairwise orthogonal imaginary simple roots (including the empty set).
\end{proposition}


\section{CY root data: the axioms}
\label{sec:cy-root-data-axioms}

Not every generalized root datum arises from a CY3 threefold.  The CY condition constrains the lattice, root multiplicities, and denominator identity.  We axiomatize the class of \emph{CY root data} as a distinguished subfamily.

\begin{definition}[CY$_3$ root datum]
\label{def:cy3-root-datum}
A generalized root datum $\mathcal{R} = (\Lambda, (\cdot,\cdot), \Delta, \mathrm{mult}, W, \rho)$ is a \emph{CY$_3$ root datum} if it satisfies axioms (CY1)--(CY7) below.
\end{definition}

\subsection*{(CY1) Lattice structure}

The lattice $\Lambda$ admits a decomposition
\[
  \Lambda = \Lambda^{\mathrm{hyp}} \oplus \Lambda^{\mathrm{ell}}
\]
where:
\begin{itemize}
  \item $\Lambda^{\mathrm{hyp}}$ is a hyperbolic sublattice of signature $(s, 1)$ with $s \geq 2$, containing all real roots.  The restriction of $(\cdot,\cdot)$ to $\Lambda^{\mathrm{hyp}}$ is non-degenerate.
  \item $\Lambda^{\mathrm{ell}}$ is a positive semi-definite sublattice (signature $(t, 0)$, $t \geq 0$) accounting for the ``elliptic directions'' (fiber classes).
  \item The total signature satisfies $(p, q) = (s + t, 1)$ or $(s + t, 2)$.  The signature $(p,2)$ case arises when the ``second negative direction'' (from the elliptic fiber class pairing) is present; the $(p,1)$ case arises for rigid CY3s.
\end{itemize}

\begin{remark}
\label{rem:cy1-mukai}
For a CY3 $X$, the lattice $\Lambda(X)$ is extracted from the Mukai lattice $\widetilde{H}(X, \mathbb{Z}) = H^0 \oplus H^2 \oplus H^4 \oplus H^6$ with the Mukai pairing, or from $K_0(X)$ with the Euler form.  When $X = S \times E$ is a K3 product, $\Lambda^{\mathrm{hyp}} = \Lambda^{2,1}_{II}$ and $\Lambda^{\mathrm{ell}} = \Lambda^{1,1}$ (the elliptic fiber lattice), giving total signature $(3,2)$.  For toric CY3, $\Lambda$ is the lattice of the toric fan with the intersection form.
\end{remark}

\subsection*{(CY2) Real root constraints}

The real roots $\Delta^{\mathrm{re}}$ lie in $\Lambda^{\mathrm{hyp}}$ and satisfy:
\begin{enumerate}[label=(\alph*)]
  \item \textbf{Positive norm}: $(\alpha, \alpha) > 0$ for all $\alpha \in \Delta^{\mathrm{re}}$.  Specifically, $(\alpha, \alpha) = 2$ for all real simple roots (the ``simply-laced'' condition).
  \item \textbf{Finiteness}: The set of real simple roots $\Pi^{\mathrm{re}} = \{\alpha_1, \ldots, \alpha_n\}$ is finite.
  \item \textbf{Off-diagonal negativity}: $(\alpha_i, \alpha_j) \leq 0$ for $i \neq j$, and $(\alpha_i, \alpha_j) \in \{0, -1, -2\}$.
  \item \textbf{Hyperbolicity}: The Gram matrix $A = ((\alpha_i, \alpha_j))$ is of hyperbolic type --- it has signature $(n-1, 1)$ or, if $n \leq s$, the submatrix is negative semi-definite off the diagonal with determinant $\leq 0$.
\end{enumerate}

\begin{remark}
\label{rem:cy2-simply-laced}
The simply-laced condition $(\alpha, \alpha) = 2$ reflects geometry: real roots arise from $(-2)$-curves or exceptional divisors, which have self-intersection $-2$ (becoming $(\alpha, \alpha) = 2$ under the Mukai pairing sign convention).  Non-simply-laced CY root data would correspond to CY3s with non-reduced fibers or orbifold singularities.
\end{remark}

\subsection*{(CY3) Imaginary root multiplicities from a Jacobi-type form}

There exists a (weak) Jacobi form or quasi-modular form $\varphi$ of weight $0$ and appropriate index such that the imaginary root multiplicities are controlled by $\varphi$:
\[
  \mathrm{mult}(\alpha) = c_\varphi(\tfrac{1}{2}(\alpha, \alpha), \ell(\alpha))
\]
where $c_\varphi(n, l)$ are the Fourier coefficients of $\varphi$ and $\ell \colon \Lambda \to \mathbb{Z}$ is a linear form (the ``elliptic index'').  More precisely:
\begin{enumerate}[label=(\alph*)]
  \item The function $n \mapsto c_\varphi(n, l)$ has at most polynomial growth for fixed $l$.
  \item $\mathrm{mult}(\alpha) > 0$ when $(\alpha, \alpha) = 0$ (isotropic imaginary roots are even/bosonic): $\Delta^{\mathrm{im}}_0 \supseteq \{\alpha \in \Delta^{\mathrm{im}} \mid (\alpha, \alpha) = 0\}$.
  \item $\mathrm{mult}(\alpha) < 0$ is permitted only when $(\alpha, \alpha) < 0$ (negative-norm roots may be odd/fermionic): $\Delta^{\mathrm{im}}_1 \subseteq \{\alpha \in \Delta^{\mathrm{im}} \mid (\alpha, \alpha) < 0\}$.
  \item The multiplicity function is $W$-invariant: $\mathrm{mult}(w \cdot \alpha) = \mathrm{mult}(\alpha)$ for all $w \in W$.
\end{enumerate}

\begin{remark}
\label{rem:cy3-bps}
The Jacobi form $\varphi$ is the BPS counting function of the CY3.  For $K3 \times E$: $\varphi = \phi_{0,1}$ (the K3 elliptic genus).  For toric CY3: $\varphi$ encodes the DT partition function, though in this case the ``Jacobi form'' interpretation requires the refined topological vertex framework.  The sign of $\mathrm{mult}(\alpha)$ encodes the fermion number of the corresponding BPS state: positive for bosonic, negative for fermionic.
\end{remark}

\subsection*{(CY4) Weyl group and reflection structure}

The Weyl group $W$ satisfies:
\begin{enumerate}[label=(\alph*)]
  \item $W$ is a Coxeter group generated by reflections $\{s_{\alpha_i}\}_{i=1}^n$ in the real simple roots.
  \item The extended symmetry group $\widetilde{W} = W \rtimes \mathrm{Aut}(\mathcal{P})$ (where $\mathrm{Aut}(\mathcal{P})$ is the automorphism group of the fundamental polyhedron) embeds into the orthogonal group $\mathrm{O}(\Lambda^{\mathrm{hyp}})_+$.
  \item The index $[\mathrm{O}(\Lambda^{\mathrm{hyp}})_+ : \widetilde{W}]$ is finite.
\end{enumerate}

\begin{remark}
\label{rem:cy4-gritsenko}
Condition (c) is the Gritsenko--Nikulin ``finite-volume'' condition: the fundamental polyhedron has finite volume in the hyperbolic space $\mathbb{H}^{s-1} \hookrightarrow C(\Lambda^{\mathrm{hyp}})_+$.  For CY3 root data this is automatic when the number of real simple roots equals the rank of $\Lambda^{\mathrm{hyp}}$, which holds in all known examples.
\end{remark}

\subsection*{(CY5) Automorphic denominator identity}

The denominator identity of $\mathcal{R}$ produces an automorphic form $\Phi_\mathcal{R}$ for $\mathrm{O}(\Lambda)_+$:
\begin{enumerate}[label=(\alph*)]
  \item The product
  \[
    \Phi_\mathcal{R}(z) = e^{-2\pi i (\rho, z)} \prod_{\alpha \in \Delta_+} (1 - e^{-2\pi i (\alpha, z)})^{\mathrm{mult}(\alpha)}
  \]
  converges for $z$ in a tube domain $\mathcal{T}_\mathcal{R} \subset \Lambda \otimes \mathbb{C}$ and extends to an automorphic form of weight $k_\mathcal{R}$ on the orthogonal symmetric domain $\mathbb{D}(\Lambda)$.

  \item The weight $k_\mathcal{R}$ satisfies
  \[
    k_\mathcal{R} = \tfrac{1}{2} c_\varphi(0,0)
  \]
  where $c_\varphi(0,0)$ is the constant term of the Jacobi form $\varphi$ from (CY3).

  \item The automorphic form $\Phi_\mathcal{R}$ is a cusp form (vanishes at all cusps) or, more generally, a modular form of ``singular weight'' in the sense of Borcherds.
\end{enumerate}

\begin{remark}
\label{rem:cy5-igusa}
For $K3 \times E$: $\Phi_\mathcal{R} = \frac{1}{64}\Delta_5$ (the Igusa cusp form).
The K3 elliptic genus in the DVV convention is
$\phi_{0,1} = 2y^{-1} + 20 + 2y + \cdots$, so the constant
term is $c_{\phi_{0,1}}(0,0) = 20$.  By axiom (CY5)(b), the
Borcherds product weight is $k_{\mathrm{Borch}} = \frac{1}{2}
\cdot 20 = 10$; the identification with $\Delta_5$ of weight $5$
uses the rescaling $\mathrm{Borch}(\phi_{0,1})(Z) =
\frac{1}{64}\Delta_5(2Z)$, so $k_\mathcal{R} = 5$.
For toric CY3: $\Phi_\mathcal{R}$ is the DT partition function,
and $k_\mathcal{R}$ is the ``Euler weight''
$\frac{1}{2}\chi(X)$.
\end{remark}

\subsection*{(CY6) Serre-type relations from the Gram matrix}

For any two simple roots $\alpha_i, \alpha_j \in \Pi^{\mathrm{re}}$ with $(\alpha_i, \alpha_j) = -m_{ij}$ (where $m_{ij} \geq 0$), the generators $e_{\alpha_i}, e_{\alpha_j}$ of the BKM superalgebra $\mathfrak{g}_\mathcal{R}$ satisfy the Serre relation
\[
  (\mathrm{ad}\, e_{\alpha_i})^{m_{ij} + 1}(e_{\alpha_j}) = 0.
\]
For a real simple root $\alpha_i$ and an imaginary simple root $\beta_j \in \Pi^{\mathrm{im}}$, the relation is:
\[
  (\mathrm{ad}\, e_{\alpha_i})^{1 - 2(\alpha_i, \beta_j)/(\alpha_i, \alpha_i)}(e_{\beta_j}) = 0
\]
when $(\alpha_i, \alpha_i) > 0$.  For two imaginary simple roots $\beta_i, \beta_j$ with $(\beta_i, \beta_j) = 0$: $[e_{\beta_i}, e_{\beta_j}] = 0$.

\subsection*{(CY7) CY integrality}

The following integrality conditions hold:
\begin{enumerate}[label=(\alph*)]
  \item $\Lambda$ is an even lattice: $(\alpha, \alpha) \in 2\mathbb{Z}$ for all $\alpha \in \Lambda$.
  \item The discriminant group $\Lambda^*/\Lambda$ is finite.
  \item The multiplicity function takes values in $\mathbb{Z}$, and the generating function $\sum_\alpha \mathrm{mult}(\alpha) \, q^{(\alpha,\alpha)/2}$ is a modular object (Jacobi form, mock modular form, or higher-depth modular form).
  \item $\rho \in \Lambda \otimes \mathbb{Q}$ with $(\rho, \rho) \in \mathbb{Q}$ (rationality of the Weyl norm).
\end{enumerate}

\begin{remark}
\label{rem:cy7-even}
The even lattice condition (a) is the CY manifestation of the absence of half-integer spin BPS states in the spectrum.  For K3 geometries this is the standard condition on the Mukai lattice; for general CY3s it follows from the integrality of the Euler form on $K$-theory.
\end{remark}


\section{Distinguishing CY root data from general BKM root data}
\label{sec:cy-vs-bkm}

The axioms (CY1)--(CY7) carve out a proper subclass of generalized root data.  We collect the expected distinguishing features into a single conjecture; a complete proof requires the full automorphic product theory.

\begin{conjecture}[Characterization of CY root data]
\label{conj:cy-characterization}
\ClaimStatusConjectured
A generalized root datum $\mathcal{R}$ is a CY$_3$ root datum if and only if the following conditions all hold:
\begin{enumerate}[label=(\arabic*)]
  \item \textbf{Lattice type}: $\Lambda$ is even, of signature $(p,1)$ or $(p,2)$ with $p \geq 2$, and decomposes as $\Lambda^{\mathrm{hyp}} \oplus \Lambda^{\mathrm{ell}}$.

  \item \textbf{Simply-laced real roots}: $(\alpha, \alpha) = 2$ for all $\alpha \in \Delta^{\mathrm{re}}$.

  \item \textbf{Modular multiplicities}: $\mathrm{mult}$ is governed by a Jacobi form $\varphi$ of weight $0$, in the sense that multiplicities depend only on the norm $(\alpha, \alpha)/2$ and the elliptic index $\ell(\alpha)$.

  \item \textbf{Bosonic isotropic / fermionic negative-norm}: $\mathrm{mult}(\alpha) > 0$ when $(\alpha, \alpha) = 0$; $\mathrm{mult}(\alpha) < 0$ only when $(\alpha, \alpha) < 0$.

  \item \textbf{Automorphic product}: The denominator product converges to an automorphic form for $\mathrm{O}(\Lambda)_+$.

  \item \textbf{Finite covolume}: $[\mathrm{O}(\Lambda^{\mathrm{hyp}})_+ : \widetilde{W}] < \infty$.

  \item \textbf{CY Euler constraint}: The weight $k_\mathcal{R}$ satisfies $k_\mathcal{R} = \frac{1}{2}\chi(X)$ for the CY3 Euler characteristic, equivalently $k_\mathcal{R} = \frac{1}{2}c_\varphi(0,0)$ with the appropriate normalization.
\end{enumerate}
\end{conjecture}

\begin{remark}[What the CY axioms exclude]
\label{rem:cy-exclusions}
Several phenomena present in general BKM root data are forbidden by the CY axioms:
\begin{enumerate}[label=(\alph*)]
  \item \emph{Non-simply-laced real roots}: General BKM algebras allow $(\alpha, \alpha) \in \{2, 4, 6, \ldots\}$.  The CY condition forces $(\alpha, \alpha) = 2$ from the geometry of $(-2)$-classes.
  \item \emph{Positive-norm fermionic roots}: In a general BKM superalgebra, odd roots can have any norm.  The CY condition requires $\Delta^{\mathrm{im}}_1 \subset \{\alpha \mid (\alpha,\alpha) < 0\}$, reflecting the spin-statistics relation for BPS states.
  \item \emph{Non-modular multiplicities}: A general BKM algebra can have arbitrary multiplicities.  The CY condition demands that the multiplicity function is the Fourier expansion of a modular object.
  \item \emph{Non-arithmetic lattices}: CY root data always have arithmetic lattices (commensurable with $\mathrm{O}(\Lambda)(\mathbb{Z})$), excluding the exotic non-arithmetic hyperbolic reflection groups.
\end{enumerate}
\end{remark}


\section{Morphisms of root data}
\label{sec:morphisms-root-data}

\begin{definition}[Morphism of generalized root data]
\label{def:morphism-root-data}
Let $\mathcal{R} = (\Lambda, (\cdot,\cdot), \Delta, \mathrm{mult}, W, \rho)$ and $\mathcal{R}' = (\Lambda', (\cdot,\cdot)', \Delta', \mathrm{mult}', W', \rho')$ be generalized root data.  A \emph{morphism} $f \colon \mathcal{R} \to \mathcal{R}'$ is a lattice homomorphism $f \colon \Lambda \to \Lambda'$ satisfying:
\begin{enumerate}[label=(\roman*)]
  \item \textbf{Isometric embedding}: $(f(\alpha), f(\beta))' = (\alpha, \beta)$ for all $\alpha, \beta \in \Lambda$.
  \item \textbf{Root preservation}: $f(\Delta^{\mathrm{re}}) \subseteq \Delta'^{\mathrm{re}}$, $f(\Delta^{\mathrm{im}}_0) \subseteq \Delta'^{\mathrm{im}}_0$, $f(\Delta^{\mathrm{im}}_1) \subseteq \Delta'^{\mathrm{im}}_1$.
  \item \textbf{Multiplicity compatibility}: $\mathrm{mult}'(f(\alpha)) \geq \mathrm{mult}(\alpha)$ for all $\alpha \in \Delta$, with equality when $f$ is an isomorphism.
  \item \textbf{Weyl equivariance}: There exists a group homomorphism $\bar{f} \colon W \to W'$ such that $f \circ w = \bar{f}(w) \circ f$ for all $w \in W$.
  \item \textbf{Weyl vector compatibility}: $f(\rho) = \pi_{\mathrm{im}(f)}(\rho')$, where $\pi_{\mathrm{im}(f)}$ denotes the orthogonal projection of $\rho'$ onto $\mathrm{im}(f) \otimes \mathbb{Q} \subset \Lambda' \otimes \mathbb{Q}$.
\end{enumerate}
\end{definition}

\begin{definition}[Isomorphism of root data]
\label{def:iso-root-data}
An \emph{isomorphism} $f \colon \mathcal{R} \xrightarrow{\sim} \mathcal{R}'$ is a morphism where $f \colon \Lambda \to \Lambda'$ is a lattice isomorphism, $f(\Delta) = \Delta'$, $\mathrm{mult}' \circ f = \mathrm{mult}$, $\bar{f}$ is a group isomorphism, and $f(\rho) = \rho'$.
\end{definition}

\begin{definition}[Sub-root-datum]
\label{def:sub-root-datum}
A \emph{sub-root-datum} $\mathcal{R}_0 \hookrightarrow \mathcal{R}$ is a morphism where $f$ is a primitive lattice embedding (i.e., $\Lambda'/f(\Lambda)$ is torsion-free) and the induced map on Weyl groups is injective.
\end{definition}

\begin{proposition}[Category of CY root data]
\label{prop:cy-root-category}
\ClaimStatusProvedHere
CY$_3$ root data with the morphisms of Definition~\ref{def:morphism-root-data} form a category $\mathbf{CYRoot}_3$.  Isomorphisms in this category are precisely the isomorphisms of Definition~\ref{def:iso-root-data}.
\end{proposition}

\begin{proof}
We verify the axioms of a category.  The identity morphism $\mathrm{id}_\mathcal{R}$ is the identity map on $\Lambda$; conditions (i)--(v) are immediate.  Composition: given $f \colon \mathcal{R} \to \mathcal{R}'$ and $g \colon \mathcal{R}' \to \mathcal{R}''$, the composite $g \circ f$ is isometric (since both are), preserves root decompositions and multiplicities (since both do), and is Weyl-equivariant via $\overline{g \circ f} = \bar{g} \circ \bar{f}$.  Weyl vector compatibility: $g(f(\rho)) = g(\pi_{\mathrm{im}(f)}(\rho')) = \pi_{\mathrm{im}(g \circ f)}(\rho'')$ by the functoriality of orthogonal projection for nested isometric images.  Associativity is inherited from the associativity of lattice homomorphisms.

It remains to check that the CY axioms are preserved: a morphism $f \colon \mathcal{R} \to \mathcal{R}'$ between CY$_3$ root data is automatically a morphism of root data (no extra condition is needed since the CY axioms are properties of the individual root data, not of the morphisms).
\end{proof}


\section{Direct sums of root data}
\label{sec:direct-sums}

\begin{definition}[Direct sum]
\label{def:direct-sum-root-data}
Given generalized root data $\mathcal{R}_1 = (\Lambda_1, (\cdot,\cdot)_1, \Delta_1, \mathrm{mult}_1, W_1, \rho_1)$ and $\mathcal{R}_2 = (\Lambda_2, (\cdot,\cdot)_2, \Delta_2, \mathrm{mult}_2, W_2, \rho_2)$, their \emph{direct sum} is
\[
  \mathcal{R}_1 \oplus \mathcal{R}_2 = (\Lambda_1 \oplus \Lambda_2,\; (\cdot,\cdot)_1 \perp (\cdot,\cdot)_2,\; \Delta_1 \sqcup \Delta_2,\; \mathrm{mult}_1 \sqcup \mathrm{mult}_2,\; W_1 \times W_2,\; \rho_1 + \rho_2)
\]
where $(\cdot,\cdot)_1 \perp (\cdot,\cdot)_2$ denotes the orthogonal direct sum of bilinear forms, and we identify $\Delta_i \subset \Lambda_i \hookrightarrow \Lambda_1 \oplus \Lambda_2$.
\end{definition}

\begin{lemma}
\label{lem:direct-sum-roots}
\ClaimStatusProvedHere
The root system of $\mathcal{R}_1 \oplus \mathcal{R}_2$ contains no ``mixed'' roots: every root in $\Delta_1 \sqcup \Delta_2$ lies entirely in $\Lambda_1$ or $\Lambda_2$.  In particular, $(\alpha, \beta) = 0$ for $\alpha \in \Delta_1$, $\beta \in \Delta_2$, and the BKM superalgebra decomposes as $\mathfrak{g}_{\mathcal{R}_1 \oplus \mathcal{R}_2} \simeq \mathfrak{g}_{\mathcal{R}_1} \oplus \mathfrak{g}_{\mathcal{R}_2}$.
\end{lemma}

\begin{proof}
By orthogonality of $\Lambda_1$ and $\Lambda_2$ in the direct sum, any vector $\alpha = \alpha_1 + \alpha_2$ with $\alpha_i \in \Lambda_i$ satisfies $(\alpha, \alpha) = (\alpha_1, \alpha_1)_1 + (\alpha_2, \alpha_2)_2$.  A real root requires $(\alpha, \alpha) = 2$; since both lattices are even (in the CY case), $(\alpha_i, \alpha_i) \in 2\mathbb{Z}$, so $(\alpha, \alpha) = 2$ forces one component to be zero.  The same argument applies to imaginary roots: the multiplicity function depends on the norm and index, which decouple in the direct sum.  The BKM algebra decomposition follows from the orthogonality of root spaces.
\end{proof}

\begin{proposition}[Direct sum and automorphic forms]
\label{prop:direct-sum-automorphic}
\ClaimStatusProvedHere
If $\mathcal{R}_1$ and $\mathcal{R}_2$ are CY root data with automorphic denominator products $\Phi_{\mathcal{R}_1}$ and $\Phi_{\mathcal{R}_2}$, then
\[
  \Phi_{\mathcal{R}_1 \oplus \mathcal{R}_2}(z_1, z_2) = \Phi_{\mathcal{R}_1}(z_1) \cdot \Phi_{\mathcal{R}_2}(z_2)
\]
where $z_i \in \Lambda_i \otimes \mathbb{C}$.  The weight is additive: $k_{\mathcal{R}_1 \oplus \mathcal{R}_2} = k_{\mathcal{R}_1} + k_{\mathcal{R}_2}$.
\end{proposition}

\begin{proof}
The product formula factorizes since $\Delta_+(\mathcal{R}_1 \oplus \mathcal{R}_2) = \Delta_+(\mathcal{R}_1) \sqcup \Delta_+(\mathcal{R}_2)$ and the exponential factors separate.
\end{proof}

\begin{warning}
\label{warn:direct-sum-not-cy}
The direct sum of two CY$_3$ root data is not, in general, a CY$_3$ root datum: the total lattice $\Lambda_1 \oplus \Lambda_2$ typically has signature $(p_1 + p_2, q_1 + q_2)$ which may violate axiom (CY1).  The direct sum is best understood as an operation on generalized root data that specializes to a CY root datum only in specific geometric situations (e.g., when $X_1 \sqcup X_2$ makes sense as a disconnected CY3, which is degenerate).  The geometrically meaningful operation for combining CY data is the fibration construction of \S\ref{sec:fibration}.
\end{warning}


\section{The fibration operation: CY$_2$ root datum $+$ elliptic curve $\to$ CY$_3$ root datum}
\label{sec:fibration}

The most important structural operation on CY root data is the \emph{fibration construction}: starting from the root datum of a CY$_2$ surface (K3 or abelian surface) and the modular data of an elliptic curve, one produces a CY$_3$ root datum via the Borcherds multiplicative lift.  This is the root-datum incarnation of the geometric fibration $X = (S \times E)/G$.

\subsection{CY$_2$ root data}
\label{ssec:cy2-root-data}

\begin{definition}[CY$_2$ root datum]
\label{def:cy2-root-datum}
A \emph{CY$_2$ root datum} is a tuple $\mathcal{R}_S = (\Lambda_S, (\cdot,\cdot)_S, \varphi_S)$ where:
\begin{enumerate}[label=(\roman*)]
  \item $\Lambda_S$ is an even lattice of signature $(r_+, r_-)$ with $r_+ \geq 1$;
  \item $(\cdot, \cdot)_S$ is the intersection form on $\Lambda_S$;
  \item $\varphi_S(\tau, z)$ is a weak Jacobi form of weight $0$ and index $1$ for $\mathrm{SL}_2(\mathbb{Z})$ (or a congruence subgroup), encoding the elliptic genus of the CY$_2$ surface $S$.
\end{enumerate}
\end{definition}

\begin{example}[K3 root datum]
\label{ex:k3-root-datum}
For a K3 surface $S$: $\Lambda_S = H^2(S, \mathbb{Z}) \simeq U^{\oplus 3} \oplus E_8(-1)^{\oplus 2}$, signature $(3, 19)$, and $\varphi_S = \phi_{0,1}$ (the K3 elliptic genus).  The Fourier expansion is
\[
  \phi_{0,1}(\tau, z) = \frac{\theta_1(\tau, z)^2}{\eta(\tau)^6} \cdot E_4(\tau) = (y^{-1} - 2 + y)(1 + \sum_{n \geq 1} c(n) q^n)
\]
with $c(1) = 20$, $c(2) = 216$, and so forth.  (Here $y = e^{2\pi i z}$, $q = e^{2\pi i \tau}$.)
\end{example}

\begin{example}[Abelian surface root datum]
\label{ex:abelian-root-datum}
For an abelian surface $A$: $\Lambda_A = H^2(A, \mathbb{Z}) \simeq U^{\oplus 3}$, signature $(3, 3)$, and $\varphi_A = \phi_{0,1}^A = -2\phi_{0,1}^{E \times E}$, which is a different Jacobi form (the abelian surface elliptic genus vanishes, so this case requires the ``refined'' version).
\end{example}

\subsection{Elliptic modular data}
\label{ssec:elliptic-data}

\begin{definition}[Elliptic datum]
\label{def:elliptic-datum}
An \emph{elliptic datum} is a pair $\mathcal{E} = (\Lambda_E, N)$ where:
\begin{enumerate}[label=(\roman*)]
  \item $\Lambda_E = \mathbb{Z} \cdot e \simeq [0]$ is a rank-1 lattice with $(e, e) = 0$ (the class of the elliptic fiber);
  \item $N \in \mathbb{Z}_{>0}$ is the ``level'' (the order of the torsion in $E[N]$ used in the quotient construction).
\end{enumerate}
\end{definition}

\subsection{The Borcherds lift}
\label{ssec:borcherds-lift}

The Borcherds multiplicative (or exponential) lift is the key mechanism.  Given a Jacobi form $\varphi$ of weight $0$ and index $1$, Borcherds's theorem produces an automorphic product on an orthogonal symmetric domain of signature $(n+2, 2)$ from a lattice of signature $(n, 0)$.  For the fibration construction, the input is the CY$_2$ root datum and the output is the CY$_3$ root datum.

\begin{construction}[Fibration of root data]
\label{constr:fibration}
Given a CY$_2$ root datum $\mathcal{R}_S = (\Lambda_S, (\cdot,\cdot)_S, \varphi_S)$ and an elliptic datum $\mathcal{E} = (\Lambda_E, N)$, the \emph{fibration root datum} $\mathcal{R}_{S,E} = \mathrm{Fib}(\mathcal{R}_S, \mathcal{E})$ is a CY$_3$ root datum constructed as follows.

\medskip\noindent
\textbf{Step 1: The CY$_3$ lattice.}  Choose a hyperbolic sublattice $\Lambda^{\mathrm{hyp}} \subset \Lambda_S$ of signature $(s, 1)$ with $s \geq 2$ (for K3, this is a sublattice of the N\'eron--Severi lattice compatible with the elliptic fibration).  Define
\[
  \Lambda = \Lambda^{\mathrm{hyp}} \oplus \Lambda^{1,1}
\]
where $\Lambda^{1,1}$ is the rank-2 hyperbolic lattice with Gram matrix $\bigl(\begin{smallmatrix} 0 & 1 \\ 1 & 0 \end{smallmatrix}\bigr)$, contributed by the fiber class and the base class of the elliptic fibration on the product $S \times E$.  The total signature is $(s+1, 2)$.

\medskip\noindent
\textbf{Step 2: Real roots.}  The real simple roots $\Pi^{\mathrm{re}} = \{\alpha_1, \ldots, \alpha_n\}$ are the primitive vectors of norm $2$ in $\Lambda^{\mathrm{hyp}}$ that bound the fundamental polyhedron $\mathcal{P}$ of the reflection group $W^{(2)}(\Lambda^{\mathrm{hyp}})$.

\medskip\noindent
\textbf{Step 3: Imaginary roots and their multiplicities.}  The imaginary roots are the vectors $\alpha \in \Lambda_+$ (in the positive cone closure) with multiplicities given by the Borcherds lift:
\[
  \mathrm{mult}(\alpha) = c_\varphi\bigl(\tfrac{1}{2}(\alpha, \alpha), \ell(\alpha)\bigr)
\]
where $c_\varphi(n, l)$ are the Fourier coefficients of $\varphi_S$ and $\ell \colon \Lambda \to \mathbb{Z}$ is the projection to the elliptic index (the $\Lambda^{1,1}$-component).  The decomposition is:
\begin{itemize}
  \item $\alpha \in \Delta^{\mathrm{im}}_0$ if $\mathrm{mult}(\alpha) > 0$ (even/bosonic);
  \item $\alpha \in \Delta^{\mathrm{im}}_1$ if $\mathrm{mult}(\alpha) < 0$ (odd/fermionic).
\end{itemize}

\medskip\noindent
\textbf{Step 4: Weyl group.}  $W = W^{(2)}(\Lambda^{\mathrm{hyp}})$, the group generated by reflections in norm-$2$ vectors of $\Lambda^{\mathrm{hyp}}$.

\medskip\noindent
\textbf{Step 5: Weyl vector.}  $\rho \in \Lambda^{\mathrm{hyp}} \otimes \mathbb{Q}$ is the unique vector in $\overline{\mathcal{P}} \otimes \mathbb{Q}$ satisfying $(\rho, \alpha_i) = -1$ for all $\alpha_i \in \Pi^{\mathrm{re}}$.
\end{construction}

\begin{theorem}[Fibration produces a CY$_3$ root datum]
\label{thm:fibration-cy3}
\ClaimStatusProvedHere
The root datum $\mathcal{R}_{S,E} = \mathrm{Fib}(\mathcal{R}_S, \mathcal{E})$ satisfies axioms (CY1)--(CY7).  The automorphic denominator product $\Phi_{\mathcal{R}_{S,E}}$ is the Borcherds multiplicative lift $\mathrm{Borch}(\varphi_S)$ of the Jacobi form $\varphi_S$.
\end{theorem}

\begin{proof}
We verify each axiom.

\emph{(CY1)}: By construction, $\Lambda = \Lambda^{\mathrm{hyp}} \oplus \Lambda^{1,1}$ with $\Lambda^{\mathrm{hyp}}$ of signature $(s,1)$ and $\Lambda^{1,1}$ of signature $(1,1)$, giving total signature $(s+1, 2)$.  Taking $\Lambda^{\mathrm{ell}} = \Lambda^{1,1}$ gives the required decomposition.

\emph{(CY2)}: The real simple roots are norm-$2$ vectors in $\Lambda^{\mathrm{hyp}}$ by Step~2.  The finiteness of $\Pi^{\mathrm{re}}$ follows from the fact that $\mathcal{P}$ is a finite-sided polyhedron (by Vinberg's algorithm applied to the hyperbolic lattice).  The off-diagonal condition $(\alpha_i, \alpha_j) \leq 0$ for $i \neq j$ holds because the $\alpha_i$ are inward-pointing normals to faces of a convex polyhedron in hyperbolic space.

\emph{(CY3)}: The multiplicity function is $\mathrm{mult}(\alpha) = c_\varphi((\alpha,\alpha)/2, \ell(\alpha))$ by definition.  Properties (a)--(d) follow from the modularity of $\varphi_S$: polynomial growth of Fourier coefficients, the isotropic roots have multiplicity $c_\varphi(0, l) > 0$ (since $\varphi_S$ is a weak Jacobi form with positive leading Fourier coefficients at $n = 0$), and $W$-invariance follows from the fact that $W$ acts trivially on $\Lambda^{1,1}$ (hence preserves $\ell$) and preserves norms (hence preserves $(\alpha,\alpha)/2$).

\emph{(CY4)}: The Weyl group is a Coxeter group by Vinberg's theory.  The finite-covolume condition holds because the fundamental polyhedron in $\mathbb{H}^{s-1}$ has finite volume (a consequence of the classification of reflective lattices; for K3 this is verified directly).

\emph{(CY5)}: This is the content of Borcherds's theorem on the multiplicative lift: the product
\[
  \Phi(z) = e^{-2\pi i (\rho, z)} \prod_{\alpha \in \Delta_+} (1 - e^{-2\pi i(\alpha, z)})^{c_\varphi((\alpha,\alpha)/2, \ell(\alpha))}
\]
converges on the tube domain over the positive cone and extends to an automorphic form for $\mathrm{O}(\Lambda)_+$ of weight $k = \frac{1}{2}c_\varphi(0,0)$.

\emph{(CY6)}: The Serre relations follow from the general theory of BKM algebras applied to the Gram matrix of $\Pi^{\mathrm{re}}$ (which has $A_{ii} = 2$ and $A_{ij} \leq 0$ for $i \neq j$ by (CY2)).

\emph{(CY7)}: $\Lambda^{\mathrm{hyp}} \oplus \Lambda^{1,1}$ is even (since $\Lambda^{\mathrm{hyp}}$ is even by the K3 lattice structure, and $\Lambda^{1,1}$ is even).  The discriminant group is finite since both summands have non-degenerate forms.  Integrality of $\mathrm{mult}$ follows from the integrality of Fourier coefficients of Jacobi forms.  Rationality of $\rho$ follows from the linear system $(\rho, \alpha_i) = -1$ with rational coefficients.
\end{proof}

\subsection{Functoriality of the fibration}
\label{ssec:fibration-functoriality}

\begin{proposition}[Functoriality of Fib]
\label{prop:fib-functorial}
\ClaimStatusProvedHere
The fibration construction defines a functor
\[
  \mathrm{Fib} \colon \mathbf{CY}_2\mathbf{Root} \times \mathbf{Ell} \longrightarrow \mathbf{CYRoot}_3
\]
where $\mathbf{CY}_2\mathbf{Root}$ is the category of CY$_2$ root data (with morphisms being isometric lattice embeddings preserving the Jacobi form) and $\mathbf{Ell}$ is the category of elliptic data (with morphisms being level-changing maps $N | N'$).
\end{proposition}

\begin{proof}
A morphism $(f, \iota) \colon (\mathcal{R}_S, \mathcal{E}) \to (\mathcal{R}_{S'}, \mathcal{E}')$ consists of a lattice embedding $f \colon \Lambda_S \hookrightarrow \Lambda_{S'}$ with $f^*\varphi_{S'} = \varphi_S$ (up to modular correction at the level change) and a divisibility $N | N'$.  The induced map on fibration root data is the lattice embedding $f \oplus \mathrm{id}_{\Lambda^{1,1}} \colon \Lambda^{\mathrm{hyp}} \oplus \Lambda^{1,1} \hookrightarrow \Lambda'^{\mathrm{hyp}} \oplus \Lambda^{1,1}$, which preserves all the CY root datum structure by the functoriality of the Borcherds lift with respect to lattice embeddings.
\end{proof}


\subsection{The K3 $\times$ E paradigm}
\label{ssec:k3-times-e-paradigm}

We spell out the fibration construction for the main example.

\begin{example}[K3 $\times$ E root datum via fibration]
\label{ex:k3-e-fibration}
Let $\mathcal{R}_S$ be the K3 root datum with $\Lambda_S = H^2(S, \mathbb{Z})$ and $\varphi_S = \phi_{0,1}$.  Choose the hyperbolic sublattice $\Lambda^{\mathrm{hyp}} = \Lambda^{2,1}_{II} \simeq \mathbb{Z}\delta_1 \oplus \mathbb{Z}\delta_2 \oplus \mathbb{Z}\delta_3$ of signature $(2,1)$ with Gram matrix
\[
  (\delta_i, \delta_j) = \begin{pmatrix} 2 & -2 & -2 \\ -2 & 2 & -2 \\ -2 & -2 & 2 \end{pmatrix}.
\]
Let $\mathcal{E} = (\mathbb{Z} \cdot e, 1)$.  Then $\mathrm{Fib}(\mathcal{R}_S, \mathcal{E})$ is the CY$_3$ root datum $\mathcal{R}(K3 \times E)$:
\begin{itemize}
  \item $\Lambda = \Lambda^{2,1}_{II} \oplus \Lambda^{1,1}$, signature $(3,2)$, as in \S\ref{sec:k3e-lattice}.
  \item $\Pi^{\mathrm{re}} = \{\delta_1, \delta_2, \delta_3\}$, $W = W^{(2)}(\Lambda^{2,1}_{II})$, $\mathrm{Aut}(\mathcal{P}) = S_3$.
  \item $\rho = f_2 - \frac{1}{2}f_3 + f_{-2}$ with $(\rho, \delta_i) = -1$ for all $i$.
  \item $\mathrm{mult}(\alpha) = f(nm, l)$ from $\phi_{0,1}$: real roots have multiplicity $1$; isotropic imaginary roots have positive multiplicity; negative-norm imaginary roots may be fermionic.
  \item $\Phi_{\mathcal{R}} = \frac{1}{64}\Delta_5$ (the Igusa cusp form of weight $5$), with $k_\mathcal{R} = 5$.
\end{itemize}
\end{example}


\section{The fibration operation at the level of algebras}
\label{sec:fibration-algebras}

The fibration of root data has a precise algebraic counterpart: it corresponds to the passage from the vertex algebra of a CY$_2$ surface to the BKM superalgebra of a CY$_3$ threefold.

\begin{construction}[From CY$_2$ vertex algebra to CY$_3$ BKM superalgebra]
\label{constr:cy2-to-cy3}
Given a CY$_2$ root datum $\mathcal{R}_S$ and an elliptic datum $\mathcal{E}$:
\begin{enumerate}[label=(\arabic*)]
  \item The CY$_2$ data determines a vertex algebra $V_S$ (the chiral algebra of the sigma model on $S$; for K3, this is related to the $\mathcal{N} = 4$ superconformal algebra at $c = 6$).
  \item The elliptic genus $\varphi_S = \mathrm{Tr}_{V_S}((-1)^F y^{J_0} q^{L_0 - c/24})$ encodes the BPS spectrum.
  \item The Borcherds lift $\mathrm{Borch}(\varphi_S) = \Phi_{\mathcal{R}_{S,E}}$ is the denominator identity of the CY$_3$ BKM superalgebra $\mathfrak{g}_{\mathcal{R}_{S,E}}$.
  \item In the language of this work: $V_S$ is the $E_1$-sector (tree-level vertex algebra); the passage $V_S \leadsto \mathfrak{g}_{\mathcal{R}_{S,E}}$ is the automorphic correction / shadow obstruction tower.
\end{enumerate}
\end{construction}

The fibration construction $\mathrm{Fib}$ is thus the root-datum shadow of the Borcherds exponential lift at the level of automorphic forms, and of the ``automorphic correction'' at the level of Lie algebras:
\[
\begin{tikzcd}[column sep=large, row sep=large]
  \mathcal{R}_S \text{ (CY$_2$)} \arrow[r, "\mathrm{Fib}(-,\mathcal{E})"] \arrow[d, "\mathfrak{g}_{(-)}"] & \mathcal{R}_{S,E} \text{ (CY$_3$)} \arrow[d, "\mathfrak{g}_{(-)}"] \\
  V_S \text{ (vertex algebra)} \arrow[r, "\text{Borcherds lift}"] & \mathfrak{g}_{\mathcal{R}_{S,E}} \text{ (BKM superalgebra)}
\end{tikzcd}
\]


\section{When does a root datum come from a CY3?}
\label{sec:realization}

The central \emph{realization problem} asks: given an abstract CY$_3$ root datum $\mathcal{R}$, does there exist a CY3 threefold $X$ with $\mathcal{R}(X) \simeq \mathcal{R}$?  We formulate necessary conditions and a partial converse.

\begin{definition}[Geometric CY root datum]
\label{def:geometric-cy-root-datum}
A CY$_3$ root datum $\mathcal{R}$ is \emph{geometric} if there exists a smooth projective CY3 threefold $X$ (or a smooth quasi-projective CY3 with at worst Gorenstein singularities) such that $\mathcal{R}(X) \simeq \mathcal{R}$ via the extraction map of \S\ref{sec:cy3-root-data}.
\end{definition}

\begin{proposition}[Necessary conditions for geometricity]
\label{prop:necessary-geometric}
\ClaimStatusProvedHere
If $\mathcal{R}$ is geometric, then:
\begin{enumerate}[label=(\roman*)]
  \item \textbf{Hodge-theoretic constraint}: The rank of $\Lambda$ is bounded by $\mathrm{rk}(\Lambda) \leq h^{1,1}(X) + h^{2,1}(X) + 2$ (from the Hodge diamond of a CY3).
  \item \textbf{Mirror constraint}: There exists a ``mirror'' CY$_3$ root datum $\mathcal{R}^\vee$ with $\Lambda^{\vee, \mathrm{hyp}}$ and $\Lambda^{\vee, \mathrm{ell}}$ exchanged (reflecting the interchange $h^{1,1} \leftrightarrow h^{2,1}$ under mirror symmetry).
  \item \textbf{Euler characteristic}: $k_\mathcal{R} = \frac{1}{2}|\chi(X)|$ where $\chi(X) = 2(h^{1,1} - h^{2,1})$ (when $h^{2,1}(X) = 0$, i.e., rigid CY3s, the signature is $(p,1)$ and $k_\mathcal{R} = h^{1,1} + 1$).
  \item \textbf{DT integrality}: The multiplicity function $\mathrm{mult}(\alpha)$ must agree with the Donaldson--Thomas invariants $\mathrm{DT}_\alpha(X)$ (which are integers satisfying the Kontsevich--Soibelman wall-crossing formula).
\end{enumerate}
\end{proposition}

\begin{proof}
(i) follows from $\mathrm{rk}(\Lambda(X)) \leq \mathrm{rk}(K_0(X)) = \sum_p (-1)^p \mathrm{rk}(K_0^p) \leq h^{1,1} + h^{2,1} + 2$ (the rank of the Mukai lattice restricted to the relevant sublattice).  (ii) is a consequence of homological mirror symmetry: if $X^\vee$ is a mirror CY3, then $D^b(\Coh(X)) \simeq D^\pi(\Fuk(X^\vee))$, and the root datum of $X^\vee$ has the exchanged Hodge numbers.  (iii) follows from the weight formula $k_\mathcal{R} = \frac{1}{2}c_\varphi(0,0)$ combined with the relation $c_\varphi(0,0) = \chi(X)$ (the constant term of the BPS generating function equals the topological Euler characteristic, up to normalization).  (iv) is the identification $\mathrm{mult}(\alpha) = \mathrm{DT}_\alpha(X)$ for a geometric CY3.
\end{proof}

\begin{conjecture}[Partial converse: fibered CY root data are geometric]
\label{conj:fibered-geometric}
\ClaimStatusConjectured
Every CY$_3$ root datum $\mathcal{R}$ that arises via the fibration construction $\mathcal{R} = \mathrm{Fib}(\mathcal{R}_S, \mathcal{E})$ from a geometric CY$_2$ root datum $\mathcal{R}_S$ is geometric.  In particular, for every reflective hyperbolic sublattice $\Lambda^{\mathrm{hyp}} \subset H^2(S, \mathbb{Z})$ with $S$ a K3 surface, the CY$_3$ root datum $\mathrm{Fib}(\mathcal{R}_S|_{\Lambda^{\mathrm{hyp}}}, \mathcal{E})$ is realized by $X = (S \times E)/G$ for an appropriate finite group $G$ acting on $S \times E$.
\end{conjecture}

\begin{remark}
\label{rem:non-fibered}
The conjecture leaves open the realizability of CY$_3$ root data that do not arise from the fibration construction.  For toric CY3s, the root datum comes from the toric fan rather than from a fibration, and realizability is automatic (the toric variety exists by construction).  The truly non-trivial cases are compact CY3s with no elliptic fibration structure, such as the quintic threefold in $\mathbb{P}^4$; for these, the root datum (if it exists in the BKM sense) is not well understood and may require an extension of the axioms.
\end{remark}


\section{Compatibility with the shadow obstruction tower}
\label{sec:shadow-tower}

We now connect the root datum axiomatics to the shadow obstruction tower of Volume~I.

\begin{conjecture}[Root datum from the shadow obstruction tower]
\label{conj:root-datum-shadow}
\ClaimStatusConjectured
Let $A_X = \Phi(\cC_X)$ be the quantum chiral algebra of a CY3 $X$ (via Theorem CY-A), and let $\Theta_{A_X} = \sum_{r \geq 2} \Theta^{(r)}_{A_X}$ be its universal MC element (shadow obstruction tower).  The CY$_3$ root datum $\mathcal{R}(X)$ is recovered as follows:
\begin{enumerate}[label=(\roman*)]
  \item The lattice $\Lambda(X)$ is the grading lattice of the bar complex $B(A_X)$.
  \item The Weyl vector $\rho$ is determined by $\Theta^{(2)}_{A_X} = \kappa(A_X)$: the arity-$2$ shadow component.
  \item The real roots correspond to the poles of the collision $r$-matrix $r(z) = \mathrm{Res}^{\mathrm{coll}}_{0,2}(\Theta^{(2)})$.
  \item The imaginary root of charge $\alpha$ with $|\alpha| = r$ (total arity) is detected by $\Theta^{(r)}_{A_X}$: the arity-$r$ shadow.  Its multiplicity is the dimension of the corresponding component of the cyclic cohomology $\HC^*_\alpha(A_X)$.
  \item The denominator identity $\Phi_{\mathcal{R}(X)}$ is the bar-complex Euler product $\chi(B(A_X))$.
\end{enumerate}
\end{conjecture}

\begin{remark}
\label{rem:shadow-dictionary}
This proposition makes the ``automorphic correction = shadow obstruction tower'' identification precise at the level of root data.  The real roots are the tree-level data (arity 2, the classical $r$-matrix).  The imaginary roots are the quantum corrections (arity $\geq 3$, the higher shadows).  The full BKM superalgebra $\mathfrak{g}_{\mathcal{R}(X)}$ is the ``infinite shadow depth'' limit of the shadow obstruction tower --- the complete automorphic correction.
\end{remark}


\section{The toric case: root data from quivers}
\label{sec:toric-root-data-quiver}

For toric CY3s, the root datum has a direct combinatorial description in terms of the quiver.

\begin{construction}[Root datum from a quiver with potential]
\label{constr:quiver-root-datum}
Let $(Q, W)$ be a quiver with potential coming from a toric CY3 $X_\Sigma$.  The CY$_3$ root datum $\mathcal{R}(Q, W)$ is:
\begin{enumerate}[label=(\roman*)]
  \item $\Lambda = \mathbb{Z}^{Q_0}$ (the lattice of dimension vectors), with bilinear form $(\mathbf{d}, \mathbf{d}') = \chi_Q(\mathbf{d}, \mathbf{d}') + \chi_Q(\mathbf{d}', \mathbf{d})$ (the symmetrized Euler--Ringel form of $Q$).
  \item $\Delta^{\mathrm{re}}$: the simple roots $\mathbf{e}_i$ (standard basis vectors) for vertices $i$ such that $(\mathbf{e}_i, \mathbf{e}_i) = 2$ (i.e., vertices with one loop; the CY condition on the potential ensures $(\mathbf{e}_i, \mathbf{e}_i) \leq 2$ always).
  \item $\Delta^{\mathrm{im}}$: dimension vectors $\mathbf{d}$ with nonzero DT invariant $\mathrm{DT}_\mathbf{d}(Q, W) \neq 0$, decomposed into even/odd by the sign of $\mathrm{DT}_\mathbf{d}$.
  \item $\mathrm{mult}(\mathbf{d}) = \mathrm{DT}_\mathbf{d}(Q, W)$ (the motivic or numerical DT invariant).
  \item $W = \langle s_i \mid i \in Q_0,\, (\mathbf{e}_i, \mathbf{e}_i) = 2 \rangle$ (reflections in the real simple roots).
  \item $\rho$: determined by $(\rho, \mathbf{e}_i) = -1$ for all real simple roots $\mathbf{e}_i$.
\end{enumerate}
\end{construction}

\begin{example}[$\mathbb{C}^3$: the Jordan quiver]
\label{ex:c3-root-datum}
$Q$ has one vertex and one loop.  $\Lambda = \mathbb{Z}$, $(\cdot, \cdot)$ is degenerate (self-intersection $= 0$ for the Jordan quiver with CY potential).  The root datum is:
\begin{itemize}
  \item No real roots ($\Delta^{\mathrm{re}} = \emptyset$).
  \item All roots are imaginary: $\Delta^{\mathrm{im}}_0 = \{n \in \mathbb{Z}_{>0}\}$ with $\mathrm{mult}(n) = p(n)$ (the partition function).
  \item $W = \{1\}$ (trivial Weyl group).
  \item The associated algebra is $Y(\widehat{\mathfrak{gl}}_1) \simeq \mathcal{W}_{1+\infty}$.
\end{itemize}
This is a degenerate CY$_3$ root datum (no real roots, degenerate form), illustrating that the axioms (CY1)--(CY7) should be understood as holding ``generically'' --- toric CY3s with no compact divisors sit at the boundary of the moduli.
\end{example}

\begin{example}[Resolved conifold: the Klebanov--Witten quiver]
\label{ex:conifold-root-datum}
$Q$ has two vertices $\{1, 2\}$ with $a$ arrows $1 \to 2$ and $b$ arrows $2 \to 1$ (for the conifold: $a = b = 2$).  The symmetrized Euler form gives a $2 \times 2$ Gram matrix.  The root datum has:
\begin{itemize}
  \item Real simple roots $\mathbf{e}_1, \mathbf{e}_2$ with $(\mathbf{e}_i, \mathbf{e}_j) = \bigl(\begin{smallmatrix} 2 & -2 \\ -2 & 2 \end{smallmatrix}\bigr)$ (the $\widehat{A}_1$ Cartan matrix).
  \item $W = \mathbb{Z}/2\mathbb{Z} \ltimes \mathbb{Z}$ (the affine Weyl group of type $\widehat{A}_1$).
  \item Imaginary roots indexed by $(n, m) \in \mathbb{Z}^2_{\geq 0}$ with DT invariants as multiplicities.
\end{itemize}
\end{example}


\section{Open problems}
\label{sec:root-datum-open}

We conclude with a list of open problems directly arising from the axiomatics.

\begin{enumerate}[label=\textbf{P\arabic*.}]
  \item \textbf{Classification of CY$_3$ root data.}  Classify all CY$_3$ root data of rank $\leq 6$.  For the fibration class ($\mathcal{R} = \mathrm{Fib}(\mathcal{R}_S, \mathcal{E})$), this reduces to classifying reflective hyperbolic lattices of rank $\leq 4$ inside K3 lattices, which is essentially the Gritsenko--Nikulin classification.  \emph{Are there CY$_3$ root data that do not arise from fibration?}

  \item \textbf{The quintic root datum.}  Construct the CY$_3$ root datum of the quintic threefold $X_5 \subset \mathbb{P}^4$.  The quintic has $h^{1,1} = 1$, $h^{2,1} = 101$, $\chi = -200$.  The lattice $\Lambda$ should be extracted from $K_0(X_5)$ with the Mukai pairing.  \emph{Is the multiplicity function governed by a (mock) Jacobi form?  If so, of what type?}

  \item \textbf{Mirror symmetry of root data.}  Formalize the mirror involution $\mathcal{R} \mapsto \mathcal{R}^\vee$ at the level of root data.  For fibered CY3s, this should correspond to the exchange of the two $\Lambda^{1,1}$ factors.  \emph{Does $\mathfrak{g}_\mathcal{R} \simeq \mathfrak{g}_{\mathcal{R}^\vee}$ (possibly up to a spectral flow)?}

  \item \textbf{Wall-crossing as Weyl chamber crossing.}  The Kontsevich--Soibelman wall-crossing formula changes the BPS spectrum as stability parameters vary.  \emph{Is wall-crossing encoded as a change of fundamental polyhedron $\mathcal{P} \mapsto \mathcal{P}'$ within the same Weyl group $W$?}  If so, gauge equivalence of MC elements (Volume~I) should correspond to the wall-crossing automorphism.

  \item \textbf{CY root data and vertex algebras.}  The CY$_3$ root datum determines a BKM superalgebra $\mathfrak{g}_\mathcal{R}$; its affinization determines a vertex algebra $V(\widehat{\mathfrak{g}}_\mathcal{R})$.  \emph{Is $V(\widehat{\mathfrak{g}}_\mathcal{R}) \simeq A_X$ (the quantum chiral algebra)?}  More precisely, is the vacuum module of the affinization isomorphic to the chiral algebra produced by the CY-to-chiral functor $\Phi$?

  \item \textbf{Higher CY root data.}  Extend the axioms to CY$_d$ root data for $d > 3$.  For CY$_4$: the automorphic form should be a Jacobi form for the Siegel modular group of genus $d - 1$; the fibration should be CY$_3$ root datum $+$ elliptic curve $\to$ CY$_4$ root datum.  \emph{What plays the role of the Borcherds lift in dimension $4$?}
\end{enumerate}
